\subsection{Noise}\label{sec:theory_general_noise}
Any unwanted signal that is present in a circuit is called noise. It is
usually measured in terms of its power spectral density (PSD) and then
converted to root mean square (RMS) voltage or current. Some common noises
are:
\begin{itemize}
  \item Thermal Noise: This is the noise generated by the thermal fluctuations
    of charge carriers (usually electrons) in a conductor or semiconductor.
    It is also called Johnson-Nyquist noise, is given by:
    \[
      S_n(f) = 4 k T R \hspace{1cm} \text{[in V\(^2\)/Hz; in PSD]} \\
      V_n = \sqrt{4 k T R \Delta f} \hspace{1cm} \text{[in V\(_{\text{rms}}\)]}
    \]
    where \(S_n(f)\) is the power spectral density of the noise, \(k\) is
    Boltzmann's constant, \(T\) is the absolute temperature in Kelvin, \(
    R\) is the resistance in ohms, and \(\Delta f\) is the bandwidth of the
    signal.
  \item Shot Noise: This is the noise generated by the discrete and probabilistic
    nature of charge carriers (usually electrons) in a conductor or semiconductor.
    It is given by:
    \[
      S_n(f) = 2 q I \hspace{1cm} \text{[in A\(^2\)/Hz; in PSD]} \\
      I_n = \sqrt{2 q I \Delta f} \hspace{1cm} \text{[in A\(_{\text{rms}}\)]}
    \]
    where \(S_n(f)\) is the power spectral density of the noise, \(q\) is
    the charge of an electron, \(I\) is the average current in amperes,
    and \(\Delta f\) is the bandwidth of the signal.
  \item Flicker Noise: This is the noise that has a power spectral density that
    is inversely proportional to frequency. It is also called 1/f noise and
    is given by:
    \[
      S_n(f) = \frac{K}{f^\alpha} \hspace{1cm} \text{[in V\(^2\)/Hz; in PSD]} \\
      V_n = \sqrt{\frac{K}{f^\alpha} \Delta f} \hspace{1cm} \text{[in V\(_{\text{rms}}\)]}
    \]
    where \(S_n(f)\) is the power spectral density of the noise, \(K\) is
    a constant that depends on the device and its operating conditions, \(f\)
    is the frequency in hertz, and \(\alpha\) is a constant that typically
    ranges from 0.5 to 2.
\end{itemize}
